\documentclass[../main.tex]{subfiles}
\begin{document}

\chapter{パイプライン}
\lessoninfo{
  video={Lesson 3，動画 1--18},
  textbook={H\&P \cite{hennessy2017} 付録 C.1--C.3},
  keywords={パイプライン，ステージ，パイプライン CPI，ストール，フラッシュ，制御依存，データ依存，RAW，WAR，WAW，ハザード，フォワーディング，パイプライン段数},
}

パイプライン化はスループットを高めるための最初にして最も基本的な技法で
ある．1つの命令を終えてから次を始めるのではなく，プロセッサは複数の命令
の実行を，それぞれ異なる進行段階で重ね合わせる．この章では，パイプライン
化されたプロセッサの動作を復習し，理想的な性能を定量化した後，理想の達成
を妨げる命令間の\emph{依存}を調べる．

\section{パイプライン化の考え方}

命令の実行は，異なるハードウェアを使う一連のステップからなる．ステップを
ラッチで区切られた\term[すてーじ]{ステージ}[stage]として配置すれば，命令
$I_1$ がステージ 1 を離れてステージ 2 に進んだ瞬間に，命令 $I_2$ が
ステージ 1 に入れる．$k$ ステージなら，最大 $k$ 個の命令が同時に処理中
となる．\source{Lesson 3，動画 2--3；H\&P §C.1．}

第2章の2つの指標に対する帰結は異なる．
\begin{itemize}
  \item \textbf{スループット}は最大 $k$ 倍に向上する．定常状態では，
    $k$ ステージ時間ごとではなく 1 ステージ時間ごとに 1 命令が完了する．
  \item 個々の命令の\textbf{レイテンシ}は改善\emph{しない}．総作業量は
    同じであり，ステージ間のラッチがわずかなオーバーヘッドを加えるため，
    命令当たりのレイテンシはむしろ少し増える．
\end{itemize}

\subsection{プロセッサにおけるパイプライン}

古典的な 5 段パイプライン（\cref{fig:stages}）を，本講義を通じて基準と
なる設計として用いる．\margin{実際のプロセッサは 10〜20 段だが，5 段モデル
でここで重要な問題はすべて捉えられる．}

\begin{figure}[htbp]
  \centering
  \input{figures/pipeline-stages}
  \caption{5 段の命令パイプライン．}
  \label{fig:stages}
\end{figure}

\begin{description}
  \item[IF] \emph{命令フェッチ}: プログラムカウンタ（PC）が指す命令を命令
    メモリから読み，PC を進める．
  \item[ID] \emph{命令デコード}とレジスタ読み出し: 演算の種類を決定し，
    ソースオペランドをレジスタファイルから読む．
  \item[EX] \emph{実行}: ALU 演算を行う，あるいはロード/ストアのメモリ
    アドレスを計算する，あるいは分岐条件を評価する．
  \item[MEM] \emph{メモリ}: ロード/ストアのためにデータメモリにアクセス
    する．\margin{各ステージはメモリアクセスも含めてちょうど 1 サイクルと
    仮定する．この仮定は第12〜15章で覆される．}
  \item[WB] \emph{書き戻し}: 結果をレジスタファイルに書く．
\end{description}

\cref{fig:timing} は独立な 5 命令がパイプラインを流れる様子を示す．最初の
4 サイクルでパイプラインが満たされた後は，毎サイクル 1 命令が完了する．
\caution{ここでは命令メモリとデータメモリが別ポートであると仮定している．
1 ポートなら\emph{構造}ハザードが加わる（H\&P 付録 C）．}

\begin{figure}[htbp]
  \centering
  % Five instructions flowing through a five-stage pipeline.
\begin{tikzpicture}[x=0.95cm,y=0.55cm, font=\footnotesize\sffamily,
  st/.style={draw, minimum width=0.9cm, minimum height=0.5cm, inner sep=0pt, fill=ln-box}]
  \foreach \c in {1,...,9} { \node at (\c,0.7) {\c}; }
  \node[anchor=east] at (0.5,0.7) {\iflangja{サイクル}{cycle}};
  \foreach \i/\name in {0/I$_1$,1/I$_2$,2/I$_3$,3/I$_4$,4/I$_5$} {
    \node[anchor=east] at (0.5,-\i) {\name};
    \foreach \s/\lab in {0/IF,1/ID,2/EX,3/MEM,4/WB} {
      \node[st] at (\i+\s+1,-\i) {\lab};
    }
  }
\end{tikzpicture}

  \caption{5 段パイプラインの命令タイミング．サイクル 5 以降は毎サイクル
    1 命令が完了する．}
  \label{fig:timing}
\end{figure}

\section{パイプライン CPI}

\cref{fig:timing} の理想的な場合，パイプラインは毎サイクル 1 命令を完了
するので，性能の鉄則（\cref{eq:iron-law}）の CPI は 1 である．
\source{Lesson 3，動画 6；H\&P §C.1．}これを 1 より大きくする効果が
2つある．
\begin{itemize}
  \item \emph{充填と排出}．最初の命令が完了するまでに $k$ サイクルかかる．
    $n$ 命令のプログラムでは合計 $k + (n-1)$ サイクル，すなわち
    CPI $= (k+n-1)/n \to 1$（$n$ が大きいとき）となる．実プログラム
    （$n$ は数十億）では無視できる．
  \item \emph{ストール}．パイプラインが命令を次のステージへ進められない
    とき，命令が完了しないままサイクルが過ぎる．パイプライン化された
    プロセッサの CPI を実際に決めるのはこのサイクルである．
    \begin{equation}
      \text{CPI} \;=\; 1 + \frac{\text{ストールサイクル数}}{\text{命令}}
      \label{eq:pipeline-cpi}
    \end{equation}
\end{itemize}

\section{ストールとフラッシュ}

パイプラインでサイクルを失う機構は2つある．\source{Lesson 3，動画 7--8；
H\&P §C.2．}

\begin{definition}[ストール]
  \term[すとーる]{ストール}[stall]とは，先行する命令は進み続ける一方で，
  ある命令（とその後続すべて）を現在のステージに 1 サイクル以上留める
  ことである．その代わりに，バブル---空きスロット---がパイプラインを
  流れていく．
\end{definition}

\begin{definition}[フラッシュ]
  \term[ふらっしゅ]{フラッシュ}[flush]とは，すでにパイプラインに
  フェッチされた命令を，そもそも実行すべきでなかったという理由で破棄する
  ことである．典型的には，結果を誤って予測した分岐の後続命令が対象と
  なる．
\end{definition}

$c$ サイクルのストールは $c$ サイクルを，$m$ 命令のフラッシュは $m$
サイクルを失わせる．フラッシュされた各命令は，有用な命令が完了できたはず
のスロットを占めていたからである．

\section{制御依存}

\term[せいぎょいぞん]{制御依存}[control dependence]とは，ある命令が
そもそも実行されるかどうかが，先行する分岐の結果に依存することである．
\source{Lesson 3，動画 9--10；H\&P §C.2．}5 段パイプラインでは，分岐の
条件と飛び先は EX の終わりまで分からない．\sidenote{初期の RISC パイプ
ラインは，無駄になるスロットをフラッシュせずに露出させた．MIPS と SPARC
では分岐の直後の命令（\emph{遅延スロット}）が必ず実行され，コンパイラが
それを埋める．処理中の命令が多くなるとスロット 1 つの効果は小さく，MIPS
Release~6（2014）はこれを廃止した．}その時点で後続の 2 命令は
すでにフェッチされている（\cref{fig:flush}）．分岐が成立していれば，
それらをフラッシュし，飛び先からフェッチをやり直さなければならない．

\begin{figure}[htbp]
  \centering
  % A branch resolved in EX: the two instructions fetched after it are flushed.
\begin{tikzpicture}[x=0.95cm,y=0.55cm, font=\footnotesize\sffamily,
  st/.style={draw, minimum width=0.9cm, minimum height=0.5cm, inner sep=0pt, fill=ln-box},
  bad/.style={st, fill=ln-caution!15, text=ln-caution}]
  \foreach \c in {1,...,8} { \node at (\c,0.7) {\c}; }
  \node[anchor=east] at (0.5,0.7) {\iflangja{サイクル}{cycle}};
  \node[anchor=east] at (0.5,0) {BEQ};
  \foreach \s/\lab in {0/IF,1/ID,2/EX,3/MEM,4/WB} { \node[st] at (\s+1,0) {\lab}; }
  \node[anchor=east] at (0.5,-1) {I$_{\text{next}}$};
  \node[bad] at (2,-1) {IF}; \node[bad] at (3,-1) {ID};
  \node[anchor=east] at (0.5,-2) {I$_{\text{next}+1}$};
  \node[bad] at (3,-2) {IF};
  \node[anchor=east] at (0.5,-3) {\iflangja{分岐先}{target}};
  \foreach \s/\lab in {0/IF,1/ID,2/EX,3/MEM,4/WB} { \node[st] at (\s+4,-3) {\lab}; }
  %% label goes below-LEFT of the arrow tip: centred under it, the (wider)
  %% Japanese word runs into the first box of the branch-target row at x=4.
  \draw[->,ln-caution] (3,0.35) -- (3,-2.4)
    node[below left, inner sep=1pt, font=\scriptsize] {\iflangja{フラッシュ}{flush}};
\end{tikzpicture}

  \caption{EX で解決される成立分岐は，2つのフェッチスロットを無駄にする．}
  \label{fig:flush}
\end{figure}

これが CPI に与えるコストは，どれだけ頻繁に起こるかに依存する．
\margin{下の \SI{20}{\percent}/\SI{60}{\percent} は講義用に丸めた数字で
ある．整数コードは実際におよそ 5 命令に 1 回分岐する（H\&P 第3章）．}講義で
使われる典型的な数値は，命令の \SI{20}{\percent} が分岐で，分岐の
\SI{60}{\percent} が成立するというものである．常に次の命令をフェッチする
単純な方針（「不成立予測」）では，成立した分岐だけがフラッシュを引き
起こす．
\begin{equation}
  \text{CPI} = 1 + 0.2 \times 0.6 \times 2 = 1.24
  \label{eq:branch-cpi}
\end{equation}

\begin{example}
  分岐がステージ 10 で解決されるより深いパイプラインでは，成立分岐ごとに
  9 命令がフラッシュされる．同じ分岐統計で CPI は
  $1 + 0.2\times 0.6\times 9 = 2.08$ となり，プロセッサは 2 サイクルに
  1 命令も完了できない．パイプラインが深いほど分岐予測（第4章）が重要に
  なるのはこのためである．
\end{example}

\tip{制御依存で失うサイクル $=$（分岐命令の割合）$\times$（予測を外す
割合）$\times$（予測ミス当たりのフラッシュ命令数）．最後の因子は
（分岐が解決するステージ）$-$ 1 である．}

\section{データ依存}

\term[でーたいぞん]{データ依存}[data dependence]は，2つの命令が同じ
レジスタまたはメモリ位置にアクセスし，少なくとも一方がそれを書き込む
ときに存在する．\source{Lesson 3，動画 11--12；H\&P §C.2．}プログラム順
でのアクセスの順序にちなんで名付けられた 3 種類がある．

\begin{description}
  \item[RAW（read after write）]%
    \index{RAW いぞん@RAW 依存}%
    後の命令が，先の命令が\emph{書く}値を
    \emph{読む}．\emph{真の依存}あるいは\emph{フロー依存}とも呼ばれる．
    データが実際に一方から他方へ流れなければならない．
  \item[WAR（write after read）]%
    \index{WAR いぞん@WAR 依存}%
    後の命令が，先の命令が\emph{読む}位置に
    \emph{書く}．\emph{逆依存}．
  \item[WAW（write after write）]%
    \index{WAW いぞん@WAW 依存}%
    両命令が同じ位置に\emph{書く}．後の
    書き込みが残らなければならない．\emph{出力依存}．
\end{description}

WAR と WAW は\term[ぎのいぞん]{偽の依存}[false dependence]あるいは
\emph{名前依存}と呼ばれる．命令間でデータは移動せず，単に同じレジスタ名
を再利用しているだけだからである．これらはリネーミング（第6章）で除去
できる．RAW 依存は除去できない．\margin{RAR---同じレジスタの 2 回の
読み出し---は依存ではない．順序は問題にならない．}

\begin{example}\leavevmode % 見出しを独立した行に置き，その下にリストを置く
  \begin{lstlisting}[language={[x86masm]Assembler}, numbers=none]
  I1:  ADD R1, R2, R3     ; R1 <- R2 + R3
  I2:  SUB R4, R1, R5     ; I1 と R1 について RAW
  I3:  MUL R1, R6, R7     ; I1 と R1 について WAW，I2 と WAR
\end{lstlisting}
\end{example}

\section{依存とハザード}

依存は\emph{プログラム}の性質であり，それが問題を起こすかどうかは
\emph{パイプライン}に依存する．\source{Lesson 3，動画 13--14；H\&P §C.2．}

\begin{definition}[ハザード]
  \term[はざーど]{ハザード}[hazard]とは，特定のパイプラインがそのまま
  進行すると依存に違反してしまう状況---依存する命令が古い値を読んだり，
  実行すべきでないときに実行されたりする状況---である．
\end{definition}

5 段パイプラインでは:
\begin{itemize}
  \item \textbf{RAW} 依存は，命令が十分近いときだけハザードになる．
    レジスタは WB で書かれ ID で読まれるので，ある命令は直前の命令と
    その前の命令に対してハザードを持つが，3 つ以上後ろの命令は書き込み後
    にレジスタを読む．\margin{サイクルの前半に書き，後半に読むレジスタ
    ファイルなら，この距離は 1 つ縮む．}
  \item \textbf{WAR} と \textbf{WAW} 依存はここではハザードにならない．
    命令は\emph{プログラム順に} ID でレジスタを読み WB で書くので，後の
    命令が先の命令を追い越すことはないからである．これらはアウトオブ
    オーダのパイプライン（第6章）でハザードになる．
  \item \textbf{制御}依存は常にハザードである．パイプラインは分岐が解決
    する前に後続命令をフェッチするからである．
\end{itemize}

\begin{keypoint}
  依存はプログラムが決め，ハザードはパイプラインが決める．同じプログラム
  でも，より深い，あるいはアウトオブオーダのパイプラインではハザードが
  増える．
\end{keypoint}

\section{ハザードの処理}

ハザードを検出したら，パイプラインには 3 つの選択肢がある．
\source{Lesson 3，動画 15--16；H\&P §C.2．}

\begin{enumerate}
  \item 依存する命令を\textbf{フラッシュ}する---誤った命令がフェッチ
    された制御ハザードの場合．
  \item 必要な値が書き込まれるまで依存する命令を\textbf{ストール}する
    ---常に正しいが，サイクルを失う．
  \item \term[ふぉわーでぃんぐ]{フォワーディング}[forwarding]（バイパス）
    で値を\textbf{補う}: 結果をレジスタファイルに到達するのを待たずに，
    生成したステージから必要とするステージへ直接送る
    （\cref{fig:forwarding}）．\margin{第9章ではストアロードフォワー
    ディングが加わる．ロードストアキューの中で，待機中のストアが後続の
    ロードに値を渡す．}
\end{enumerate}

\begin{figure}[htbp]
  \centering
  % RAW dependence resolved by forwarding from EX/MEM to EX.
\begin{tikzpicture}[x=0.95cm,y=0.9cm, font=\footnotesize\sffamily,
  st/.style={draw, minimum width=0.9cm, minimum height=0.5cm, inner sep=0pt, fill=ln-box},
  >={Stealth[length=1.6mm]}]
  \foreach \c in {1,...,6} { \node at (\c,0.7) {\c}; }
  \node[anchor=east] at (0.5,0) {\texttt{ADD R1,R2,R3}};
  \foreach \s/\lab in {0/IF,1/ID,2/EX,3/MEM,4/WB} { \node[st] (a\s) at (\s+1,0) {\lab}; }
  \node[anchor=east] at (0.5,-1) {\texttt{SUB R4,R1,R5}};
  \foreach \s/\lab in {0/IF,1/ID,2/EX,3/MEM,4/WB} { \node[st] (b\s) at (\s+2,-1) {\lab}; }
  \draw[->,thick,ln-tip] (a2.south east) to[out=-60,in=120] (b2.north west);
  %% The label sits in the gap between the two rows, to the right of the arc.
  %% Centred at (4.6,-0.35) the (wider) Japanese word prints over the MEM and
  %% WB boxes of the ADD row; y=0.9 opens a gap tall enough to hold it.
  \node[ln-tip, font=\scriptsize, anchor=west] at (3.75,-0.5)
    {\iflangja{フォワーディング}{forwarding}};
\end{tikzpicture}

  \caption{フォワーディング: \texttt{ADD} の ALU 結果を \texttt{SUB} の
    EX ステージへ直接送り，ストールを回避する．}
  \label{fig:forwarding}
\end{figure}

フォワーディングは RAW ストールの大半を除去するが，すべてではない．値を
\emph{時間を遡って}送ることはできないからである．\tip{ストールサイクル数
$=$（値が用意できるサイクル）$-$（値が必要なサイクル）．$\le 0$ なら
フォワーディングだけで足りる．}ロードの結果は MEM の
終わりにならないと得られないが，次の命令はそれを自身の EX の開始時---同じ
サイクル---に必要とする．したがって\emph{ロード-使用}の組では，完全な
フォワーディングがあっても 1 サイクルのストールが避けられない．

\begin{example}\leavevmode % 見出しを独立した行に置き，その下にリストを置く
  \begin{lstlisting}[language={[x86masm]Assembler}, numbers=none]
  LW  R1, 0(R2)      ; R1 は MEM の終わり（サイクル 4）に得られる
  ADD R3, R1, R4     ; EX（サイクル 4）で R1 が必要: 1 ストール
\end{lstlisting}
  フォワーディングありでストールなしだと \texttt{ADD} は古い R1 を使って
  しまう．フォワーディングなしなら 2 サイクルストールする．
  \margin{Skylake のコアでは L1 ヒットに 4 サイクルかかるため，ロード-使用
  の組のストールは 1 ではなく約 3 サイクルになる（Intel 最適化マニュアル）．}
\end{example}

\section{ステージ数はいくつにすべきか}

$k$ 段にパイプライン化するとスループットが最大 $k$ 倍になるなら，なぜ
数百段にしないのか．\source{Lesson 3，動画 17；H\&P §C.1，§1.5．}深さを
増やすと 3 つの効果が異なる方向に働く．
\begin{itemize}
  \item \textbf{クロック周期が短くなる}（ステージ当たりの論理が減る）．
    これは性能の鉄則の $T_{\mathrm{clk}}$ を下げる，意図された利点である．
    しかしステージ当たりのラッチのオーバーヘッドは縮まないので，利得は
    頭打ちになる．\caution{パイプラインを深くしても必ず速くなるとは限ら
    ない．決めるのは積 CPI $\times T_{\mathrm{clk}}$ である（第2章）．}
  \item \textbf{CPI が増える}: ステージが増えると，分岐とその解決の間に
    処理中の命令が増え，予測ミス当たりのフラッシュ命令数が増える．また
    RAW ハザードにストールが必要な距離も伸びる．\margin{分岐予測ミス
    1 回のコストは Skylake 級のコアで約 16〜20 サイクルである
    （Agner Fog，マイクロアーキテクチャマニュアル）．}
  \item \textbf{電力が増える}: 毎サイクル切り替わるラッチが増え，周波数の
    上昇が動的電力（\cref{eq:active-power}）を押し上げる．
\end{itemize}

性能のみ（CPI $\times T_{\mathrm{clk}}$ の積）を考えると最適は深く，
30〜40 段程度である．電力を考慮に入れると最適はおよそ 10〜15 段に移り，
これが現在の設計が位置する範囲である．\sidenote{x86 で出荷された最も深い
パイプラインは 2004 年の Pentium~4「Prescott」で，31 段，最高
\SI{3.8}{\giga\hertz} であり，電力がそこで頭打ちにした．後継の Core
マイクロアーキテクチャ（2006）は 14 段に戻り，以後の深さは 15〜20 段前後
にとどまっている．}

\begin{keypoint}[章のまとめ]
  パイプライン化はスループットを上げるがレイテンシは上げない．理想 CPI
  は 1 で，すべてのストールとフラッシュがそれに加算される
  （\cref{eq:pipeline-cpi}）．制御依存と RAW 依存はサイクルを失わせる
  ハザードとなる．フォワーディングは RAW ストールの大半を除去し，分岐予測
  （次章）が制御ハザードに対処する．パイプラインの深さはクロック周期と
  CPI・電力とのトレードオフである．
\end{keypoint}

\end{document}
